% TODO make hw stylesheet
\documentclass[10pt]{article}
\usepackage[letterpaper,top=1in]{geometry}
\usepackage[dvipsnames,svgnames]{xcolor}
\usepackage[shortlabels]{enumitem}
\usepackage[framemethod=TikZ]{mdframed}
\usepackage{amsmath,amssymb,amsthm}
\usepackage[colorlinks]{hyperref}
\usepackage{multicol}
\usepackage{tabularx}

\hypersetup{
	colorlinks=true,
	linkcolor=blue,
	filecolor=magenta,
	urlcolor=magenta,
	pdftitle={MAT 528 HW}
}

\usepackage{mathtools}
\usepackage{thmtools}
\usepackage[textsize=footnotesize]{todonotes}
\usepackage{titling}

\reversemarginpar
\mdfdefinestyle{mdbluebox}{roundcorner=10pt,innerbottommargin=9pt,
	linecolor=blue,backgroundcolor=TealBlue!5,}
\declaretheoremstyle[headfont=\sffamily\bfseries\color{MidnightBlue},
mdframed={style=mdbluebox},]{thmbluebox}

\mdfdefinestyle{mdbloobox}{frametitlefont=\bfseries,innerbottommargin=8pt,
	nobreak=true,backgroundcolor=TealBlue!5,linecolor=blue,}
\declaretheoremstyle[headfont=\bfseries\color{MidnightBlue},
mdframed={style=mdbloobox},headpunct={\\[3pt]},postheadspace=0pt,]{thmbloobox}

\mdfdefinestyle{mdredbox}{frametitlefont=\bfseries,innerbottommargin=8pt,
	nobreak=true,backgroundcolor=Salmon!5,linecolor=RawSienna,}
\declaretheoremstyle[headfont=\bfseries\color{RawSienna},
mdframed={style=mdredbox},headpunct={\\[3pt]},postheadspace=0pt,]{thmredbox}

\mdfdefinestyle{mdgreenbox}{linecolor=ForestGreen,backgroundcolor=ForestGreen!5,
	linewidth=2pt,rightline=false,leftline=true,topline=false,bottomline=false,}
\declaretheoremstyle[headfont=\bfseries\sffamily\color{ForestGreen!70!black},
mdframed={style=mdgreenbox},headpunct={ --- },]{thmgreenbox}

\mdfdefinestyle{mdorangebox}{linecolor=RedOrange,backgroundcolor=RedOrange!5,
	linewidth=2pt,rightline=false,leftline=true,topline=false,bottomline=false,}
\declaretheoremstyle[headfont=\bfseries\sffamily\color{RedOrange!70!black},
mdframed={style=mdorangebox},headpunct={ --- },]{thmorangebox}

\mdfdefinestyle{mdblackbox}{linecolor=black,backgroundcolor=RedViolet!5!gray!5,
	linewidth=3pt,rightline=false,leftline=true,topline=false,bottomline=false,}
\declaretheoremstyle[mdframed={style=mdblackbox}]{thmblackbox}

\declaretheoremstyle[spaceabove=6pt,spacebelow=6pt]{basehead}

\declaretheorem[style=basehead,name=Problem]{problem}
\declaretheorem[style=thmredbox,name=Problem,numbered=no]{problem*}
\declaretheorem[style=thmbloobox,name=Setup]{setup} % for config geo
\declaretheorem[style=thmredbox,name=Example,sibling=problem]{example}
\declaretheorem[style=thmredbox,name=$\bigstar$ Problem,sibling=problem]{niceprob}
\declaretheorem[style=thmbluebox,name=Theorem,sibling=problem]{theorem}
\declaretheorem[style=thmbluebox,name=Corollary,sibling=theorem]{corollary}
\declaretheorem[style=thmbluebox,name=Proposition,sibling=theorem]{prop}
\declaretheorem[style=thmbluebox,name=Lemma,numberwithin=problem]{lemma}
\declaretheorem[style=thmbluebox,name=Theorem,numbered=no]{theorem*}
\declaretheorem[style=thmbluebox,name=Lemma,numbered=no]{lemma*}
\declaretheorem[style=thmgreenbox,name=Claim,numberwithin=problem]{claim}
\declaretheorem[style=thmgreenbox,name=Claim,numbered=no]{claim*}
\declaretheorem[style=thmblackbox,name=Remark,numberwithin=problem]{remark}
\declaretheorem[style=thmblackbox,name=Remark,numbered=no]{remark*}
\declaretheorem[style=thmgreenbox,name=Definition,sibling=theorem]{definition}
\declaretheorem[style=thmgreenbox,name=Definition,numbered=no]{definition*}

\providecommand{\ol}{\overline}
\providecommand{\ceil}[1]{\left\lceil #1 \right\rceil}
\providecommand{\floor}[1]{\left\lfloor #1 \right\rfloor}
\newcommand{\dd}{\mathrm{d}}
\providecommand{\ul}{\underline}
\providecommand{\eps}{\varepsilon}
\providecommand{\half}{\frac{1}{2}}
\providecommand{\CC}{\mathbb C}
\providecommand{\FF}{\mathbb F}
\providecommand{\NN}{\mathbb N}
\providecommand{\QQ}{\mathbb Q}
\providecommand{\RR}{\mathbb R}
\providecommand{\ZZ}{\mathbb Z}
\providecommand{\dg}{^\circ}
\providecommand{\ii}{\item}
\providecommand{\setdiff}{\smallsetminus}
\providecommand{\alert}[1]{{\sffamily\textbf{\textcolor{blue}{#1}}}}
\providecommand{\inv}{^{-1}}
\providecommand{\mb}[1]{\mathbf{#1}}

\newenvironment{soln}{\begin{proof}[Solution]}{\end{proof}}

\providecommand{\pow}{\operatorname{Pow}}
\providecommand{\vocab}[1]{{\textbf{\textcolor{ForestGreen}{#1}}}}
\providecommand{\lcm}{\operatorname{lcm}}
\providecommand{\abs}[1]{\left|#1\right|}

\usepackage{mathrsfs}

% place title at top right
\setlength{\droptitle}{-8ex}
\pretitle{\begin{flushright}\Large\bfseries}
\posttitle{\par\end{flushright}}
\preauthor{\begin{flushright}}
\postauthor{\end{flushright}}
\predate{\begin{flushright}}
\postdate{\end{flushright}}

\title{MAT 528 HW05}
\author{\large Aditya Pahuja}
\date{\large Due 10/02}
\begin{document}
\maketitle

\begin{problem}
    [Munkres 24.1]
    \begin{enumerate}[(a)]
	\ii Show that no two of the spaces $(0,1)$, $(0,1]$, and $[0,1]$
	are homeomorphic.
	\ii Suppose that there exist embeddings $f\colon X\to Y$
	and $g\colon Y\to X$. Show by means of an example
	that $X$ and $Y$ need not be homeomorphic.
	\ii Show $\RR^n$ and $\RR$ are not homeomorphic if $n > 1$.
    \end{enumerate}
\end{problem}

\begin{soln}
    (a) First, let $A$ be either of $(0,1]$ or $[0,1]$
    and suppose that $h\colon A\to (0,1)$ is a homeomorphism.
    Then, $A\setminus\{1\}$ is connected (because it is either
    $(0,1)$ or $[0,1)$), so
    \begin{equation*}
	h(A\setminus\{1\}) = h(A)\setminus\{h(1)\} = (0,1)\setminus\{h(1)\}
    \end{equation*}
    is also connected (where the first equality is because $h$ is bijective).
    However, $(0,1)\setminus\{h(1)\} = (0,h(1))\cup(h(1),1)$
    is disconnected because it is the union of nonempty disjoint open sets
    $(0,h(1))$ and $(h(1),1)$.
    Therefore, $(0,1)$ is not homeomorphic to either of $(0,1]$ or $[0,1]$.

    Now, let $g\colon[0,1]\to(0,1]$ be a homeomorphism. Then
    \begin{equation*}
	g([0,1)) = (0,1]\setminus\{g(1)\}
    \end{equation*}
    is connected, which means $g(1) = 1$, since otherwise
    this set is the union of the nonempty disjoint open sets
    $(0,g(1))\cup(g(1),1]$ (they are open in $(0,1]$). Similarly
    \begin{equation*}
	g((0,1]) = (0,1]\setminus\{g(0)\}
    \end{equation*}
    is connected, also implying $g(0) = 1$;
    this contradicts the fact that $g$ must be bijective to be a homeomorphism,
    so $(0,1]$ and $[0,1]$ are not homeomorphic.
    \\

    (b) $X = (0,1)$ and $Y = [0,1]$ are not homeomorphic, but
    $X$ can be embedded in $Y$ by the inclusion map $f(x) = x$,
    and $Y$ can be embedded in $X$ by the map $g\colon Y\to X$ defined as
    \begin{equation*}
	g(t) = \frac{1 + 2t}4;
    \end{equation*}
    this map is continuous and injective, and its inverse
    $g\inv\colon g(X)\to X$ is continuous because it is $g\inv(t) = \frac{4t-1}2$.
    \\

    (c) Let $h\colon\RR^n\to\RR$ be a homeomorphism.
    Then, $\RR^n\setminus\{0\}$ is connected because it is path connected:
    for most $x,y\in\RR^n$, the segment between $x$ and $y$ is a path
    in $\RR^n\setminus\{0\}$ connecting the two points;
    the only exception is if the segment contains $0$, in which case
    we can take the path consisting of the two segments
    between $x$ and $p$ and between $p$ and $y$, where
    $p$ is a point not on the line through $x$ and $y$.
    This means $h(\RR^n\setminus\{0\}) = \RR\setminus\{h(0)\}$ is connected,
    but this is false, because $\RR\setminus\{h(0)\}$ is the union
    of the nonempty disjoint open sets $(-\infty,h(0))\cup(h(0),\infty)$,
    so $h$ cannot exist.
\end{soln}

\pagebreak

\begin{problem}
    [Munkres 24.2]
    Let $f\colon S^1\to\RR$ be a continuous map.
    Show that there exists a point $x\in S^1$ such that $f(x) = f(-x)$.
\end{problem}

\begin{soln}
    Define $g\colon S^1\to\RR$ by $g(x) = f(x) - f(-x)$;
    we want to show that there exists $x_0\in S^1$ such that $g(x_0) = 0$.
    We can see that $g(-x) = f(-x) - f(x) = -g(x)$.
    If $g(x) = 0$ we are done; otherwise, $0$ lies between $g(x)$ and $g(-x)$
    so, since $S^1$ is connected, the intermediate value theorem tells us
    that there exists $x_0\in S^1$ such that $g(x_0) = 0$, as desired.
\end{soln}

\begin{problem}
    [Munkres 24.3]
    Let $f\colon X\to X$ be continuous.
    Show that if $X = [0,1]$, there is a point $x$ such that $f(x) = x$.
    The point $x$ is called a \vocab{fixed point} of $f$.
    What happens if $X$ equals $[0,1)$ or $(0,1)$?
\end{problem}

\begin{soln}
    If $X = [0,1]$, let $g(x) = f(x) - x$. Then $g(0) = f(0) - 0 \ge 0-0 = 0$,
    while $g(1) = f(1)-1 \le 1-1 = 0$, so $0$ is between $g(0)$ and $g(1)$;
    by the intermediate value theorem,
    there exists $x_0\in X$ such that $g(x_0) = 0$, which implies $f(x_0) = x_0$.
    \\

    If $X = (0,1)$ or $[0,1)$, then $f(x) = \half(x + 1)$ satisfies
    \begin{equation*}
        f(x) = \half x + \half \in(x,1)
    \end{equation*}
    for all $x < 1$, so this $f$ doesn't have a fixed point
    and the statement fails to hold.
\end{soln}

\begin{problem}
    [Munkres 24.11]
    If $A$ is a connected subspace of $X$, does it follow that
    $\operatorname{Int}A$ and $\partial A$ are connected?
    Does the converse hold?
\end{problem}

\begin{soln}
    Both the statement and its converse are false.

    First, $X=[0,1]$ is connected but has boundary $\{0,1\}$, which is disconnected
    (because it inherits the discrete topology as a subspace of $\RR$),
    so the boundary of a connected set may be disconnected.
    To show that the interior of a connected set may be disconnected,
    let $X\subseteq\RR^2$ be the union of the two closed disks
    of radius $1$ centered at $(-1,0)$ and $(1,0)$
    so the disks intersect only at the point $(0,0)$;
    this nonempty intersection implies that $X$ is connected
    because the disks themselves are connected.
    However, the interior of $X$ doesn't include $(0,0)$
    (and of course no other point on the $y$-axis since
    $X$ only intersects the $y$-axis at $(0,0)$)
    because any open ball around $(0,0)$ contains some point $(0,r)$ for $r\ne 0$,
    so $X$ has the separation $(X\cap H_-)\sqcup(X\cap H_+)$,
    where $H_-\subseteq\RR^2$ is the open set consisting of
    the points with a negative $x$-coordinate
    and $H_+\subseteq\RR^2$ is the open set consisting of
    the points with a positive $x$-coordinate.

    For the converse, we note that $\QQ$ has empty interior
    (since any nonempty open interval contains an irrational number ---
    the irrational numbers are dense in $\RR$)
    and its boundary is $\ol\QQ = \RR$
    (minus the interior of $\QQ$, but that's empty),
    so the interior and boundary of $\QQ$ are connected.
    However, $\QQ$ itself is disconnected,
    for example having the separation
    $(\QQ\cap(-\infty,\sqrt2))\sqcup(\QQ\cap(\sqrt2,\infty))$.
\end{soln}

\begin{problem}
    [Munkres 25.4]
    Let $X$ be locally path connected.
    Show that every connected open set in $X$ is path connected.
\end{problem}

\begin{soln}
    By Munkres's Theorem 25.3, since $X$ is locally path connected,
    each path component of an open subset $U\subseteq X$
    is open in $X$. If $U$ has multiple path components,
    then let $P$ be one of the path components
    and $Q$ be the union of the other path components;
    both $P$ and $Q$ are open in $U$, nonempty, disjoint, and have union $U$
    (because the path components of $U$ partition $U$),
    so $U$ is disconnected. Contrapositively, if $U$ is connected,
    then it must have only one path component, i.e. $U$ is also path connected.
\end{soln}

\pagebreak

\begin{problem}
    [Munkres 25.5]
    Let $X$ denote the rational points of the interval $[0,1]\times\{0\}$ of $\RR^2$.
    Let $T$ denote the union of all line segments joining the point $p=(0,1)$
    to points of $X$.
    \begin{enumerate}[(a)]
        \ii Show that $T$ is path connected,
	but is locally connected only at the point $p$.
	\ii Find a subset of $\RR^2$ that is path connected
	but is locally connected at none of its points.
    \end{enumerate}
\end{problem}

\begin{soln}
    (a) $T$ is path connected because
    given $x$ and $y$ in $T$, there is a path traversing $x\to p\to y$
    via the segments joining $x$ to $p$ and $p$ to $y$.
    By the same argument, the intersection $T\cap B_\eps(p)$
    with an open ball centered at $p$ is path connected,
    since if $x\in T\cap B_\eps(p)$ then the segment from $p$ to $x$
    is contained in $T\cap B_\eps(p)$ (if this weren't true,
    then, noting that the segment is in $T$, we would have
    a contradiction to the fact that $B_\eps(p)$ is convex).
    Thus, for any open set $U\subseteq T$ containing $p$,
    taking $\eps>0$ small enough so that $T\cap B_\eps(p)\subseteq U$
    shows that $T$ is locally path connected at $p$.

    Now consider $T^* = T\setminus \{p\}$, which is an open subset of $T$.
    Let $f\colon T^*\to\QQ\cap[0,1]$ be the map such that
    points on the segment in $T*$ connecting $(0,1)$ to a rational point $(r,0)$
    (but of course not containing $(0,1)$ because we excluded it from $T^*$)
    for $r\in\QQ\cap[0,1]$ are sent to $r$.
    This map is continuous because the pre-image of the open ball
    $(q_1,q_2)\cap\QQ$ is the union of the segments corresponding to
    the elements of that set, which make up the open subset
    \begin{equation*}
	T^*\cap\left\{(x,y): 1 - \frac{x}{q_1} < y < 1 - \frac{x}{q_2}\right\}
    \end{equation*}
    (intersecting $T^*$ with the region between the two lines).

    For each point $x$ in $T^*$, a neighborhood $U$ of that point intersects
    infinitely many of the segments that comprise $T^*$,
    so $f(U)$ is an infinite subset of $\QQ\cap[0,1]$.
    With the subspace topology, $\QQ\cap[0,1]$ has only the singleton sets
    as its connected components, so $f(U)$ is disconnected, hence so is $U$.
    Since no open subset of $x$ is connected, $T^*$, hence $T$,
    is not locally connected at $x$.
    \\

    (b) Rotate $T$ by $\pi$ about $(\half,\half)$ to get $T'$,
    with $p' = (1,0)$ being the image of $p$ under the rotation.
    The map $f$ and the corresponding map
    $f'\colon T'^*\setminus\{p'\}\to\QQ\cap[0,1]$
    given by applying $f$ to the pre-image of each point in $(T')*$
    under the rotation are both continuous, and $f,f'$ agree on $T^*\cap (T')^*$
    because they are both $1$ on this segment, so by the pasting lemma
    $g\colon T^*\cup (T')^* = T\cup T'\to \QQ\cap[0,1]$
    defined by $g|_{T^*} = f$ and $g|_{(T')^*} = f'$ is continuous.
    Then, one can apply the argument of (a) to get that
    any neighborhood of any point of $f$ is disconnected, despite
    path connectedness holding since the path connected sets $T$ and $T'$
    intersect on the segment between $p$ and $p'$ which is path connected.
\end{soln}










\end{document}

